\documentclass[12pt,a4paper,twoside]{article}
\usepackage[T1]{fontenc}
\usepackage[utf8]{inputenc}
\begin{document}
Voici un paragraphe normal (avec une
indentation de première ligne nulle pour
gagner de la place). Cet exemple suffit
à voir l’alignement des bords gauche et
droit.

\begin{flushright}
Voici un paragraphe au fer à droite où
l’alignement ne se fait plus qu’au
niveau du bord droit. La syntaxe
anglaise peut être traduite par
\emph{aligné~à~droite}.
\end{flushright}\par % pour changer un peu

\begin{flushleft}
Un paragraphe au fer à gauche. La syntaxe
anglaise est évidemment cohérente en
remplaçant droite~par~gauche.
\end{flushleft}

\begin{center}
Paragraphe centré. Seules les lignes
relativement courtes ont un intérêt sinon,
on ne voit pas très bien le caractère
centré de la chose !
\end{center}

\end{document}